\chapter{Taxonomy Construction Formalism}\label{app:dag-formalism}

This appendix gives the full architectural and mathematical detail
behind the taxonomy construction pipeline summarised in
Section~\ref{sec:arch-taxonomy} (Chapter~\ref{ch:system-architecture})
and Section~\ref{sec:meth-construction-summary}
(Chapter~\ref{ch:methodology}). The construction pipeline is a
necessary prerequisite that the arena evaluation runs on, not the
thesis's own headline contribution; this appendix exists so the exact
mechanics are preserved and traceable without requiring the reader to
work through them before reaching the arena's methodology and
results.

\section{Construction Pipeline Walkthrough}\label{sec:app-dag-pipeline}

\subsection{Seed Initialisation from Ground-Truth Domains}

Before iterative refinement begins, each corpus query is assigned to
the anchor node corresponding to its MMLU-Pro ground-truth domain label.
This creates 14 top-level partitions---one per MMLU-Pro domain---each
instantiated as a \texttt{GraphNode} with an initial member set
$Q_D^{\text{canonical}}$ equal to the set of queries carrying label
$D$.
(In the implementation the anchors sit at depth~1 beneath a single
formal root node at depth~0; the text refers to anchors as the top
level throughout, and depth-conditioned rules below are stated in
implementation depths.)

Anchor immutability is enforced structurally: every destructive
topology operation---starved-leaf pruning, sibling and redundant
fusion, and sole-child dissolution---exempts nodes at depth $\leq 1$,
so anchors are never removed, merged, or dissolved and appear in every
serialised snapshot with the same identifier.
The anchor's \texttt{sliceDim} is fixed at $d = 256$, the single
embedding dimension used at every depth
(Section~\ref{subsec:dim-schedule}).

The ground-truth label assignment serves as initialisation only.
The trickle routing phase may reassign a query to a geometrically more
appropriate anchor, so that the adapted partition
$Q_D^{\text{adapted}}$ at the end of construction need not equal
$Q_D^{\text{canonical}}$.
The distinction is recorded explicitly in the snapshot and is
material for the validation comparison: the comparison is performed
using $Q_D^{\text{adapted}}$, which reflects the geometry-corrected
membership rather than the editorial label.

\subsection{Iterative Six-Phase Refinement Loop}

After anchor initialisation, \texttt{TaxonomyEngine} executes a
six-phase loop over the whole taxonomy (the loop is global: splitting
proceeds per depth across all branches, and queries may migrate
between anchors at every iteration, Section~\ref{subsec:migration}).
The loop runs for a fixed iteration budget, with completion semantics
defined in Section~\ref{subsec:quiescence}.
Table~\ref{tab:pipeline} summarises the phases.

\begin{table}[h]
\centering
\caption{Taxonomy construction phases, implementing classes, and outputs.}
\label{tab:pipeline}
\begin{tabular}{clll}
\toprule
Phase & Responsibility & Implementing class & Output \\
\midrule
1 & MRL prefix selection and $\ell_2$-normalisation
& \texttt{EmbeddingService}
& Normalised slice matrix $X \in \mathbb{R}^{N \times 256}$ \\
2 & vMF MLE and NiW posterior fitting
& \texttt{TaxonomyFitter}
& $(\hat{\boldsymbol{\mu}}, \hat{\kappa})$ per node;
NiW hyperparameters $(m_N, \kappa_N, \nu_N, \Lambda_N)$ \\
3 & Soft routing (skipped in iteration 1)
& \texttt{TaxonomyTrickler}
& Updated weighted member sets; residual pools \\
4 & Adaptive splitting (proposal-gated)
& \texttt{TaxonomySplitter}
& New children where local gates hold and $\Delta J > 0$ \\
5 & Topological optimisation (proposal-gated)
& \texttt{TaxonomyMerger}
& Fused siblings; pruned starved leaves;
dissolved sole children \\
6 & Convergence check
& \texttt{TaxonomyEngine}
& GED quiescence streak; early stop or next iteration \\
\bottomrule
\end{tabular}
\end{table}

\textbf{Phase 1 --- Embedding selection.}
At every depth, \texttt{EmbeddingService} extracts the fixed
256-dimensional prefix slice from the pre-computed MRL vectors and
$\ell_2$-normalises each row to the unit hypersphere.

\textbf{Phase 2 --- vMF/NiW fitting.}
\texttt{TaxonomyFitter} fits a single-component vMF to each node's
member embeddings, producing the directional centroid
$(\hat{\boldsymbol{\mu}}, \hat{\kappa})$, and subsequently a
Normal-inverse-Wishart posterior retained for diagnostic scoring.
The concentration fit excludes residual-flagged queries; the mean
direction includes them.
Statistical detail is given in Section~\ref{sec:meth-vmf}.

\textbf{Phase 3 --- Trickle routing.}
\texttt{TaxonomyTrickler} routes each query down the tree under three
separated criteria---a parameter-free descent-versus-residual gate,
direction-only sibling competition with a cosine beam, and a final
per-query membership-share floor
(Section~\ref{sec:meth-routing}).
Queries that reach no destination are retained as residual mass at
the node the walk converged to, with weight and embedding preserved
(Section~\ref{subsec:residuals}).
This phase is skipped entirely in iteration~1: the ground-truth
bootstrap assignment is kept for the first split/optimise/refit pass,
and geometric reassignment starts at iteration~2
(Section~\ref{subsec:migration}).

\textbf{Phase 4 --- Adaptive splitting.}
\texttt{TaxonomySplitter} proposes partitions of non-anchor nodes via
vMF-EM clustering in a PCA-projected subspace; a proposal must pass
the local gates of Section~\ref{sec:meth-split}
(chance-corrected separation, routed sustainability, sibling
distinctness) and is then committed only if the measured global
objective improves beyond the neutrality band
($\Delta J > \tau$; Section~\ref{sec:meth-proposal-gate}).

\textbf{Phase 5 --- Topological optimisation.}
\texttt{TaxonomyMerger} executes a sequence of anchor-exempt
operations, every one of which is applied as a measured proposal
under the gate of Section~\ref{sec:meth-proposal-gate}:
starved-leaf handling (a three-way choice between keeping the leaf,
absorbing it into its parent, and fusing it into its nearest
sibling, decided by measured $J$); fusion of insufficiently separated
siblings; fusion of near-duplicate pairs that share a parent; and
dissolution of structurally redundant sole children.
Every one of these operators preserves the single-parent invariant
(Section~\ref{subsec:polyhierarchy}), so each iteration maps a tree to
a tree.

\textbf{Phase 6 --- Convergence check.}
At the end of each iteration, \texttt{TaxonomyEngine} computes the
graph edit distance to the preceding iteration's structure and either
stops construction early once the structure has been quiescent for five
consecutive iterations, triggers the next iteration, or, when the
budget is exhausted, serialises the taxonomy as an immutable snapshot
(Section~\ref{subsec:quiescence}).

\subsection{Cross-Domain Migration}\label{subsec:migration}\label{sec:meth-migration}

At bootstrap, each anchor $D_k$ is initialised with the canonical
partition---every query placed on its ground-truth anchor with weight
$1.0$---and a vMF/NiW self-model fit at $d = 256$; the anchor's
identity and top-level position are set once and never modified.

This bootstrap assignment is deliberately preserved through the first
construction iteration: Phase~3 trickle reassignment is skipped in
iteration~1, so the first split/optimise/refit pass operates on the
pure ground-truth partition, and geometric reassignment starts at
iteration~2, once real sub-domain structure exists for routing to act
on.
(Re-routing immediately in iteration~1---even with a ground-truth
bias nudge, \texttt{enableGtWarmStart}, which is disabled---let
queries drift onto merely-similar wrong anchors before any structure
existed; the iteration-1 skip removes the need for any such bias.)
From iteration~2 onward, queries are reassigned by the standard
trickle routing of Section~\ref{sec:meth-routing};
\textit{migration} is emergent from this re-routing rather than a
separately gated mechanism with its own margin, and the
\textit{adapted partition} $Q_{D_k}^{\text{adapted}}$ is defined as
the anchor-level assignment after construction completes.

Three downstream effects follow from migration.
First, the vMF self-model of each anchor is refit in Phase~2 of
subsequent iterations on $Q_{D_k}^{\text{adapted}}$, not
$Q_{D_k}^{\text{canonical}}$.
Second, all sub-domain structure induced below anchor $D_k$ is built
from $Q_{D_k}^{\text{adapted}}$, meaning that a query that migrated
into $D_k$ may reach a leaf that is thematically distant from its
MMLU-Pro label.
Third, the validation comparison (Section~\ref{sec:meth-validation})
is performed using $Q_{D_k}^{\text{adapted}}$ at the anchor level.
The migration rate
$m_{D_k} = |Q_{D_k}^{\text{adapted}} \setminus
Q_{D_k}^{\text{canonical}}| / |Q_{D_k}^{\text{canonical}}|$ is
reported as a diagnostic descriptor: it quantifies the fraction of
queries that changed anchor assignment and contextualises the
interpretation of the validation comparison.
The primary RQ2 metric is the fidelity difference
$\Delta\tau = \tau_{\text{adapted}} - \tau_{\text{canonical}}$---the
Kendall rank-agreement gain from using geometry-corrected rather than
editorial-label partitions when computing anchor-level accuracy---with
the Spearman-based $\Delta\rho$ reported secondarily.

\subsection{Completion, Quiescence, and Reproducibility}\label{subsec:quiescence}

Construction termination is criterion-based with a budget ceiling:
the loop runs for at most \texttt{numIterations} iterations (value
specified in Chapter~\ref{ch:experimental-design}), and a structural
early-stopping gate---enabled in the canonical configuration
(\texttt{enableEarlyStopping})---can terminate it earlier.
At the end of each iteration, the \texttt{TaxonomyStabilizer} counts
the structural edits between the current taxonomy $\mathcal{G}^{(t)}$
and the preceding iteration's $\mathcal{G}^{(t-1)}$: node additions and
deletions plus edge additions and deletions.
After a minimum iteration count (proportional to the number of
domains, floored at $5$), an iteration counts as quiescent when this
edit count is zero; construction stops early once five consecutive
iterations are quiescent.
Canonical runs converge at approximately iteration~15.
A leaf-level log-semantic-volume trajectory is additionally tracked as
a convergence diagnostic.
If the streak never completes, the budget is exhausted and the
snapshot is frozen as-is; downstream analyses can distinguish
quiescent from still-evolving taxonomies from the recorded GED
trajectory.

Construction is deterministic under a fixed seed: PCA power-iteration
initialisation is seeded from the problem shape, node identifiers are
monotonic counters (hash-set traversal and proposal-enumeration
orders therefore repeat exactly), and order-sensitive floating-point
aggregations iterate their inputs in sorted order.
Two runs at the same seed produce structurally identical taxonomies
(same nodes, edges, and populations), which makes cross-configuration
comparisons exact.

\subsection{Residual Retention}\label{subsec:residuals}

A query whose walk terminates at an internal node without reaching
any admissible destination is not dropped and not force-assigned: it
is retained at that node with its full unit weight, its embedding,
and a residual flag.
Three properties follow.
Mass is conserved ($\sum_v \sum_q w_v(q) = N$ holds at every
proposal, and is asserted at runtime).
Internal nodes hold hard queries \emph{iff} those queries are
residual-flagged (the C3 invariant), so residual mass is always
distinguishable from routed mass.
Residual pools are first-class objects: they are measurable (each
pool is a cell of the global objective $J$,
Section~\ref{sec:meth-proposal-gate}), routable (a pool node is an
admissible destination for later walks), recoverable (a coherent pool
can be carved into new children by the splitter,
Section~\ref{sec:meth-split}).
Residual mass is therefore an explicit measurement of
central or cross-domain content that the current structure does not
yet explain, rather than a failure mode.

\section{Directional Node Model}\label{sec:meth-vmf}

\subsection{von Mises--Fisher Density and MLE}

The vMF distribution on $\mathcal{S}^{d-1}$ with mean direction
$\boldsymbol{\mu}$ and concentration $\kappa \geq 0$ has density
\[
p(\mathbf{x} \mid \boldsymbol{\mu}, \kappa)
= C_d(\kappa) \exp\!\left(\kappa \boldsymbol{\mu}^\top \mathbf{x}\right),
\qquad
C_d(\kappa) = \frac{\kappa^{d/2 - 1}}{(2\pi)^{d/2}\, I_{d/2 - 1}(\kappa)}.
\]
Given $N$ unit-vector observations, the MLE of $\boldsymbol{\mu}$ is
the normalised sample mean $\hat{\boldsymbol{\mu}} =
\bar{\mathbf{x}} / \|\bar{\mathbf{x}}\|_2$.
The Banerjee closed-form approximation~\parencite{banerjee2005vmf}
\[
\hat{\kappa}_{0}
= \frac{\bar{R}\, d - \bar{R}^3}{1 - \bar{R}^2},
\qquad \bar{R} = \|\bar{\mathbf{x}}\|_2,
\]
avoids iterative root-finding for $\kappa$ and is accurate for
$d \geq 4$, $N \geq 5$, but is only an approximation to the true MLE
and is biased upward in high dimensions at
small sample sizes---an $O(d/N)$ bias that becomes severe when the
ratio $d/N$ is large.

Three refinements are applied, all reflecting what the
implementation actually executes.
First, the closed-form value seeds up to five Newton--Raphson steps
on the exact likelihood equation $A_d(\kappa) = \bar{R}$, where
$A_d(\kappa) = I_{d/2}(\kappa)/I_{d/2-1}(\kappa)$, using
$A_d'(\kappa) = 1 - A_d(\kappa)^2 - \tfrac{d-1}{\kappa}A_d(\kappa)$;
the seed is close enough that convergence typically occurs within one
or two steps, yielding $\hat{\kappa}_{\text{ML}}$.
Second, the refined estimate is multiplied by a small-sample
shrinkage factor
\[
\hat{\kappa}
= \hat{\kappa}_{\text{ML}} \cdot \frac{N - 1}{N + d - 2},
\]
a multiplicative correction in the spirit of the high-accuracy
$\kappa$ estimators of Hornik and Gr\"{u}n~\parencite{hornik2014vmf},
which pulls the estimate down toward its small-sample
expectation.
Because this correction scales with $N$, nodes fitted on larger
populations---parents---are systematically sharper than their
children, which is the reason routing decisions compare
\emph{directions} rather than densities
(Section~\ref{sec:meth-routing}).
Third, the shrunk estimate is blended with a parent-anchored
empirical-Bayes prior $\kappa_{\text{parent}}$ (the mean of the
node's parents' concentrations, or a default prior at the root):
\[
\kappa_{\text{final}}
= (1 - \alpha(\rho))\,\hat{\kappa}
+ \alpha(\rho)\,\kappa_{\text{parent}},
\qquad
\alpha(\rho) =
\begin{cases}
0 & \rho \leq 2,\\
(\rho - 2)/8 & 2 < \rho \leq 10,\\
1 & \rho > 10,
\end{cases}
\]
where $\rho = d/n$ is the dimension-to-effective-sample-size ratio,
so a node's own data dominate when it is well supported and the
parent prior takes over smoothly as support thins.
The joint choice of fixed dimension and minimum leaf size should be
read against this ramp honestly: at $d = 256$ and $n_{\min} = 50$
(Section~\ref{sec:exp-env}), a minimum-size leaf sits at
$\rho \approx 5.1$ and hence at a blend weight of
$\alpha \approx 0.39$---partially prior-anchored, not comfortably
outside any special regime; only nodes with $n \gtrsim 128$ are
estimated from their own data alone, and the fully prior-dominated
$\rho > 10$ regime is additionally logged with a warning.
Numerically stable evaluation of $\log I_\nu(\kappa)$ via the Debye
asymptotic expansion is given in Appendix~\ref{app:numerics}.

\subsection{NiW Auxiliary Regulariser}

A Normal-inverse-Wishart (NiW) posterior is fit to each node's member
embeddings as an auxiliary Bayesian summary of the node's local
geometry; it does not participate in routing, which is vMF-based
throughout (Section~\ref{sec:meth-routing}).
The NiW is \textit{not} the conjugate prior for the vMF---the vMF's
conjugate is the vMF itself~\parencite{gopal2014mises}---but is used
as an engineering regulariser whose posterior predictive is a
multivariate Student-$t$, accurate when $\kappa$ is moderate to large
so that the tangent-plane approximation to $\mathcal{S}^{d-1}$ holds.
Near-uniform nodes ($\kappa < 0.5$) are treated as diffuse and are
excluded from ordinary splitting unless their residual pool passes
the residual-viability gate of Section~\ref{sec:meth-split}.

The prior mean is the parent's vMF centroid, the prior pseudocount is
$\kappa_0 = 1.0$, the prior degrees of freedom are
$\nu_0 = d + 2$, and the prior scale is
$\boldsymbol{\Lambda}_0 = \lambda_0 \mathbf{I}_{d}$ with
$\lambda_0 = 1/(\hat{\kappa}_{\text{parent}} d)$.
Because the construction fixes $d = 256$ at every depth
(Section~\ref{sec:meth-mrl}), a diagonal approximation to
$\boldsymbol{\Lambda}_N$ is stored at all depths, reducing memory from
$O(d^2)$ to $O(d)$ per node and making cross-depth node comparisons
directly valid without down-projection.
The posterior update equations and the resulting Student-$t$
predictive density are given in Appendix~\ref{app:numerics}; in the
current implementation the stored posterior serves diagnostic roles
(notably the leaf log-semantic-volume convergence diagnostic of
Section~\ref{subsec:quiescence}), while all routing, migration, and
topology decisions operate on the vMF geometry of
Section~\ref{sec:meth-routing}.
Node centroids are refit directly to the fresh weighted MLE each
iteration, subject to a no-update band that freezes a node's centroid
once its refit direction agrees with the previous one to within
cosine $0.9975$; an earlier exponential-moving-average smoothing of
the centroid was removed after matched A/B runs showed it amplifies
oscillation rather than damping it.

\section{Fixed-Dimension Matryoshka Embedding Slice}\label{sec:meth-mrl}\label{subsec:dim-schedule}

MRL~\parencite{kusupati2022mrl} trains a single embedding model such
that every prefix slice $\mathbf{x}_{1:m}$ is itself a meaningful
embedding, allowing the full vector to be stored once and sliced at
query time.
Its key property is \textit{coarse-to-fine information packing}: the
most discriminative semantic directions are concentrated in the
leading coordinates, so any prefix $\mathbf{x}_{1:m'}$ is a coherent
lower-dimensional representation rather than a truncated artefact.

TaxoArena uses a single fixed slice
\[
d(\ell) = 256 \quad \text{for all depths } \ell,
\]
rather than a per-depth dimension schedule.
The choice balances two pressures.
A larger $d$ increases angular resolution and improves the
discriminability of fine-grained sub-topic splits, but the vMF
concentration estimate degrades as $d/N$ grows
(Section~\ref{sec:meth-vmf}); at $d = 256$ the fully prior-dominated
regime ($d/N > 10$) is crossed only below $N \approx
26$, and the $n_{\min} = 30$ floor keeps every fitted leaf above
that.
A minimum-size leaf nevertheless sits high on the empirical-Bayes
blend ramp: at $N = 30$, $\rho = d/N \approx 8.5$ gives
$\alpha = (\rho - 2)/8 \approx 0.82$, so its concentration estimate is
largely parent-anchored rather than fit from its own scatter.
This is a direct cost of the small-leaf setting, and it is why leaf
size is a calibrated parameter rather than a free one.
Fixing a single dimension additionally means all nodes---across all
depths---live on the same unit hypersphere, so sibling-fusion and
dissolution comparisons require no cross-depth
down-projection.
The coarse-to-fine property of MRL is exploited not by widening the
slice with depth but by the recursive structure itself: shallow nodes
separate coarse domain-level directions, while deeper splits resolve
progressively finer sub-topic structure within the same 256-dimensional
geometry.
(One qualification: split \emph{proposals} are clustered in a
PCA-projected subspace of 32, 64, or 128 dimensions depending on
cluster size, Section~\ref{sec:meth-split}; all fitting, routing, and
acceptance decisions remain at the fixed 256 dimensions.)
Re-normalisation after slicing is mandatory, since MRL guarantees
semantic meaningfulness of $\mathbf{x}_{1:d}$ but not unit norm:
\(
\mathbf{x}^{(\ell)}
= \mathbf{x}_{1:d} / \|\mathbf{x}_{1:d}\|_2.
\)
The dimension itself is ablated (A1, Section~\ref{sec:appendix-ablations})
at $\{128, 256, 512\}$ to confirm that $d = 256$ is a
quality-stable operating point rather than a fragile one.

\section{Trickle Routing}\label{sec:meth-routing}

Routing separates three decisions that a single absolute floor
previously conflated: whether to descend at all, which siblings
compete, and which reached destinations count as memberships.
At an internal node $v$, the competition set $K(v) = C(v)$ is the
node's children.

\textbf{(i) Descent-versus-residual gate.}
The walk descends from $v$ for query $\mathbf{x}$ iff
\[
\max_{c \in K(v)}
\hat{\boldsymbol{\mu}}_c^\top \mathbf{x}
\;\geq\;
(\bar{r}_v - \delta)\,
\hat{\boldsymbol{\mu}}_v^\top \mathbf{x},
\qquad
\bar{r}_v
= \Big\| \sum_{c \in C(v)} \omega_c\, \hat{\boldsymbol{\mu}}_c \Big\|_2
\in [0,1],
\]
where $\omega_c$ are the children's normalised query masses and
$\delta = \texttt{descentMargin} \geq 0$ (canonical $0$).
The factor $\bar{r}_v$ makes the bound tight: the parent's direction
is approximately the \emph{normalised} mean of its children's, so the
raw parent dot product overstates what the children can collectively
achieve by exactly the factor $1/\bar{r}_v$; after correction the
test reads ``some child must match the query at least as well as the
children's own mass-weighted average alignment.''
At $\delta = 0$ the gate is parameter-free, and a query is residual
only when it is genuinely unexplained by the sub-structure;
$\delta > 0$ relaxes the bound, trading measured residual mass for
softer leaf populations.
The comparison is between directions, never densities: the
$N$-dependent shrinkage of Section~\ref{sec:meth-vmf} makes parents
systematically sharper than children, so log-density comparisons
across levels are structurally biased toward the parent.
If the gate fails, the query is retained as residual mass at $v$
(Section~\ref{subsec:residuals}).

\textbf{(ii) Direction-only sibling competition with a cosine beam.}
Admitted siblings are scored
$f_c = \bar{\kappa}\,\hat{\boldsymbol{\mu}}_c^\top \mathbf{x}$ with
the \emph{shared mean} concentration $\bar{\kappa}$ of the level;
per-child concentrations and normalisers are excluded from the
competition, so a child cannot attract queries through concentration
bookkeeping, only through direction.
The beam keeps every sibling within an additive cosine margin of the
best,
\[
\Big\{ c :
\hat{\boldsymbol{\mu}}_c^\top \mathbf{x}
\geq
\max_{c'} \hat{\boldsymbol{\mu}}_{c'}^\top \mathbf{x} - \gamma
\Big\},
\qquad \gamma = \texttt{routingBeamGamma},
\]
the argmax always survives, and per-level transition probabilities
renormalise over the beam.
The walk accumulates log path mass; a node reached along several
paths aggregates them by log-sum-exp, so a multi-parent destination's
membership is the sum over its incoming paths.
Paths whose absolute posterior falls below $10^{-4}$ are abandoned as
a purely numerical fan-out guard carrying no membership semantics.

\textbf{(iii) Final membership share.}
After the walk, the query's accumulated masses are normalised over
the destinations it actually reached, and a destination counts as a
genuine membership iff it holds at least
\texttt{membershipFloor} of \emph{that query's own} normalised
membership.
Because the floor is self-normalised, its meaning is invariant to
depth and fan-out---unlike an absolute product-versus-floor test,
which mathematically forbids balanced structure beyond shallow depths
and forces dominant-child chains.
If no destination clears the share, the single best destination is
kept: trimming negligible tails is this floor's job, declaring
residuals is the gate's.

Membership weights over the reached destinations are normalised per
query, and the \textit{primary} leaf is the argmax of these
weights---a bookkeeping convention used for aggregation under the
primary-path convention, not an exclusivity property of the fits,
which use the full normalised weights.
Held-out queries at evaluation time are routed by exactly the same
gate, beam, and share floor, applied read-only to the frozen
snapshot, with the reached-leaf set additionally capped at
\texttt{maxLeafAssignments} per query as a judge-call-cost bound
(construction-time membership is unbounded); there is no separate
evaluation-time scoring model.
MMLU-Pro labels are never used in the routing decision; they
are inspected only afterward, for the canonical-versus-adapted
divergence metrics (Section~\ref{sec:meth-validation}).

\section{Adaptive Splitting Oracle}\label{sec:meth-split}

Split proposals are generated by vMF-EM clustering; for efficiency
the clustering runs in a PCA-projected subspace whose dimension
depends on cluster size (32 dimensions below 100 members, 64 below
500, 128 otherwise), with deterministic, seeded power-iteration
initialisation; the resulting partition is then re-fit and evaluated
in the fixed 256-dimensional routing geometry, so no acceptance
decision is made in the reduced space.
For an internal node, the split target is its residual pool---the
mechanism by which coherent unexplained mass is carved into new
children alongside existing ones.

A proposed partition $\{S_c\}$ is judged by a
\emph{chance-corrected separation score}.
For a set $S$ of unit vectors, define the total pairwise cosine
dissimilarity
\[
W(S) = |S|^2 - \Big\|\sum_{\mathbf{x} \in S} \mathbf{x}\Big\|_2^2
= \sum_{i \neq j \in S} \left(1 - \mathbf{x}_i^\top \mathbf{x}_j\right),
\]
computable in $O(|S|\,d)$ from the member sum alone.
The observed within-cluster scatter of the partition is
$\sum_c W(S_c)$; under a uniformly random partition of the same $n$
points into the same cluster sizes $\{n_c\}$, its exact expectation is
\[
\mathbb{E}\Big[\sum_c W(S_c)\Big]
= W(S_{\text{total}}) \cdot
\sum_c \frac{n_c (n_c - 1)}{n (n - 1)},
\]
because each unordered pair lands inside cluster $c$ with probability
$n_c(n_c-1)/(n(n-1))$.
The score is
\[
s = 1 - \frac{\sum_c W(S_c)}
{\mathbb{E}\left[\sum_c W(S_c)\right]},
\]
which equals $1$ for a perfect partition (each cluster internally
identical), is approximately $0$ for a random partition, negative for
an anti-clustered one, and exactly $0$ for the degenerate $k = 1$
``partition''---so a non-separating wrapper can never clear a
positive threshold.
The chance correction removes the mechanical dependence a raw
within/total ratio has on $k$ and on the cluster-size profile, so a
single threshold means the same thing everywhere in the graph.
This score is \emph{not} a Dasgupta cost delta and carries no
Dasgupta optimality interpretation; Dasgupta's cost
\parencite{dasgupta2016cost} is retained only as the whole-tree
Total Dasgupta Cost evaluation metric (Appendix~\ref{app:metrics}).

Acceptance proceeds in four stages, of which the EM clustering is
only the first.
\begin{enumerate}
\item \textbf{Proposal.} vMF-EM with $k \in \{2, \ldots, 4\}$,
$k$ selected by chance-corrected marginal improvement
$\geq \epsilon_{\text{sep}}$; every EM cluster must hold at least
$n_{\min}$ members.
\item \textbf{Routed sustainability.} Proposal vMF components are
refit at $d = 256$ and every target query is re-assigned by the same
level-local posterior the trickler uses; the split is rejected unless
\emph{every} child holds at least $n_{\min}$ members under that
routed assignment.
This closes the loop between birth and survival: a split occurs only
if routing will sustain every child it creates, which eliminated a
persistent spawn-and-starve churn in which the majority of pruned
nodes died within two iterations of birth.
Children are then refit on their routed populations.
\item \textbf{Separation and distinctness.} The routed partition
must score $s \geq \epsilon_{\text{sep}}$ ($2\epsilon_{\text{sep}}$
for small populations), and each new child must remain
chance-corrected-distinct ($\geq \epsilon_{\text{sep}}$) from every
existing sibling branch.
\item \textbf{Global veto.} The split is committed only if the
measured global objective does not degrade, $\Delta J > 0$
(Section~\ref{sec:meth-proposal-gate}); local quality is already
guaranteed by stages 1--3, so the global check is a veto rather than
an additional toll.
\end{enumerate}
The recursion depth is hard-capped at $D_{\max}$.
The threshold $\epsilon_{\text{sep}}$ (\texttt{proposalSeparationBar})
has no closed-form optimum and
trades over-splitting against under-splitting; its canonical value is
$\epsilon_{\text{sep}} = 0.02$, subject to final calibration.
% TODO(calibration): finalize once taxonomy construction is frozen
The \emph{same} score at the \emph{same} $\epsilon_{\text{sep}}$
also gates sibling fusion, sibling distinctness, and residual-split
viability, unifying these decisions under a single dimensionless
separability scale.

\section{The Proposal Gate}\label{sec:meth-proposal-gate}

Every structural edit---split, fusion, prune, and dissolution---is
applied as a \emph{measured proposal} rather than by
its local rule alone.
The gate operates on a global objective $J$: the chance-corrected
separation score of Section~\ref{sec:meth-split} lifted to the whole
corpus, with one cell per leaf (soft membership weights) plus one
cell per internal node's residual pool, so that residual mass is
measured rather than ignored.

The proposal procedure is:
snapshot the current state (topology, node
parameters, and assignments); apply the candidate edit tentatively;
recompute descent-gate statistics; clear and fully re-route every
query through the modified taxonomy; measure $J' = J(\mathcal{G}')$; and
commit under a lexicographic rule on $(J, -|V|)$, where $|V|$ is the
node count:
\[
\text{accept}
\iff
\Delta J > \tau
\quad\text{or}\quad
\bigl( |\Delta J| \leq \tau \;\wedge\; \Delta |V| < 0 \bigr),
\qquad
\Delta J = J' - J .
\]
Otherwise the snapshot is restored.
An edit is therefore taken when it strictly improves the global
objective, or when it leaves the objective unchanged and strictly
simplifies the structure.
The tolerance $\tau$ (\texttt{tau}, canonical $10^{-6}$) is a
neutrality band, not a quality bar: each edit's own quality gate---the
separation score for splits, the similarity pre-filter for
fusions---has already been applied upstream, so the gate does not
charge a second, edit-type-specific toll for the same question.
The rule is uniform across edit types, which is what makes accept and
reject decisions comparable between splits and shrink edits.

Figure~\ref{fig:proposal-gate-flow} summarises this procedure as a
single flowchart.

\begin{figure}[h]
\centering
\fbox{\parbox{0.9\textwidth}{\centering\vspace{0.6cm}
Figure E.1 --- Proposal-gate flow diagram (placeholder)\\
\footnotesize Compact flowchart: snapshot $\rightarrow$ tentative apply $\rightarrow$ no-edge early exit $\rightarrow$ re-route $\rightarrow$ lexicographic $(\Delta J, \Delta|V|)$ test $\rightarrow$ commit/restore; structural, no run data required.\vspace{0.6cm}}}
\caption{Proposal-gate flowchart. The candidate-edit path runs from
snapshot through tentative application, the no-edge early exit,
full re-routing of every query, and the lexicographic
$(\Delta J, \Delta |V|)$ test against the neutrality band $\tau$, to
commit or restore.}
\label{fig:proposal-gate-flow}
\end{figure}

Figure~\ref{fig:proposal-gate-scatter} shows the gate operating as a
measured decision boundary on the canonical run's proposals.

\begin{figure}[h]
\centering
\fbox{\parbox{0.9\textwidth}{\centering\vspace{0.6cm}
Figure E.2 --- Proposal-gate accept/reject scatter (placeholder)\\
\footnotesize Scatter of the logged proposals from the canonical seed-42 run: $\Delta |V|$ vs.\ $\Delta J$ (log-symmetric x-axis), marker shape by edit type, fill by acceptance, neutrality band $\tau$ shaded; data: \texttt{headless\_run.log} proposal lines.\vspace{0.6cm}}}
\caption{Accepted and rejected structural proposals under the gate.
Each proposal from the canonical seed-42 construction log is plotted
at its measured $\Delta J$ (symmetric log scale) and node-count change
$\Delta |V|$, with marker shape encoding the edit type (split, fusion,
prune, dissolution) and fill encoding acceptance; the neutrality band
$|\Delta J| \leq \tau$ is shaded, so the edits admitted purely for
simplifying the structure are the ones inside the band with
$\Delta |V| < 0$.}
\label{fig:proposal-gate-scatter}
\end{figure}

Two structural consequences are worth stating.
Sole-child dissolution needs no separate justification: a parent's
only tree child separates nothing (the score of
Section~\ref{sec:meth-split} is identically $0$ at $k = 1$), and the
gate confirms empirically that removing it does not degrade $J$;
top-level nodes are exempt.
Starved-leaf handling (a leaf whose branch falls below $n_{\min}$,
the same floor that gates birth) is a three-way choice---keep,
absorb into parent, or fuse into the nearest sibling---decided by
whichever option measures best under the gate, not by a fixed rule.

\section{Polyhierarchy: a Measured Negative Result}\label{subsec:polyhierarchy}\label{sec:meth-topology}\label{sec:poly-negative}

The construction produces a strict tree: every node below the root has
exactly one parent.
An earlier version of the pipeline admitted polyhierarchy through a
dedicated operator. This section records that operator, the evidence
that retired it, and the invariant that replaced it.

Two claims have to be kept apart.
Construction-time routing is weighted multi-membership
(Section~\ref{sec:meth-routing}), which lets a \emph{query} belong to
several leaves.
Polyhierarchy is a different claim: it lets a \emph{concept} belong to
several branches, through node-level topology edits.
The first is retained; the second is not.

\textbf{The operator that was tested.}
A cross-link was a growth edit: an existing node $n$ gained a second
parent $h$ when $h$'s residual pool demonstrably fitted $n$.
The motivating case is a genuinely cross-domain query---a
biostatistics question between Mathematics and Biology---which under a
strict tree can only pick one branch, fail to specialise there, and
residualise at the domain anchor.
Candidates were scored by how many of $h$'s residuals they
\emph{captured}, with capture applying the runtime descent gate itself:
\[
\hat{\boldsymbol{\mu}}_n^\top \mathbf{x}
\;\geq\;
(\bar{r}_h - \delta)\,
\hat{\boldsymbol{\mu}}_h^\top \mathbf{x}.
\]
``Captured'' therefore meant ``will actually descend at $h$ once the
edge exists,'' not proxy similarity.
Guards required acyclicity, no grandparent or redundant edges, and
disjoint top-level anchors for $h$ and $n$.
Acceptance used the proposal gate with the growth-side threshold and
the \emph{target} as the site, $\Delta J > \epsilon_{\text{sep}} \cdot
\pi_n$, so that the threshold scaled with the mass the new edge could
actually move.

\textbf{What the measurements showed.}
Across 42 cross-link proposals on the canonical corpus, the capture
fraction $f$ that the operator uses to rank candidates was
uncorrelated with the objective it must satisfy: $r(f, \Delta J) =
+0.085$, with complete distributional overlap between accepted and
rejected proposals (median $f$ of $0.802$ accepted against $0.818$
rejected).
The effect sizes were an order of magnitude apart: median $\Delta J$
was $6.9 \times 10^{-5}$ for cross-links against $5.3 \times 10^{-4}$
for splits.
Admission was therefore decided by evaluation order rather than by
evidence, which the proposal sequence makes explicit: one node was
accepted at four hosts and rejected at two, and acquired five parents
by iteration~1.

\textbf{Why a partition objective cannot decide this.}
$J$ scores a partition of query mass into cells.
Giving a node a second parent changes which paths reach its cell, but
it leaves the cell contents almost unchanged, so $\Delta J$ carries
little signal about whether the edge is semantically right.
A partition-based objective is structurally indifferent to
polyhierarchy, and no threshold on such an objective can separate good
cross-links from bad ones.

\textbf{The decision, and what it costs.}
The operator was removed rather than retuned.
The multi-membership claim survives that removal: the multi-leaf rate
of held-out queries was $0.1549$ with cross-links and $0.1518$ without,
so queries that legitimately span domains still reach several leaves
through the routing beam and the membership floor.
What is given up is the \emph{concept}-level claim, which the evidence
above could not support in the first place.

\textbf{The single-parent invariant.}
Removing the generator is not by itself sufficient, because fusion also
redirects parent edges: fusing two nodes that sit under different
parents leaves the survivor holding both, which reintroduces
multi-parent nodes with no cross-link operator present.
Fusion proposals are therefore restricted to pairs sharing a parent,
which makes every topology operator a map from a tree to a tree.
The structural-integrity gates of Appendix~\ref{app:tuning-protocol}
check the invariant directly by counting multi-parent nodes, which must
be zero.

Evaluation-time multi-leaf assignment of held-out queries---the same
routing, applied read-only to the frozen snapshot and capped at
\texttt{maxLeafAssignments}---is a separate mechanism: it does not
modify the frozen topology but
determines which leaves a given query contributes comparisons to
during the arena run (Section~\ref{sec:arch-arena}).
